\scanpage{90}
\margin{- 175 -}
\margin{IX}

(13.3.15)
\diagram{\begin{tikzcd}
0 \arrow[r] & R^f/R^t \arrow[r] & E/E^t \arrow[r] & B \arrow[r] & 0
\end{tikzcd}}
.

On notera que
\[
\text{(13.3.16)}\qquad R^f/R^t=(R_o/R_o^t)_S
\]
est un groupe analytique constant sur $S$. Il résulte donc de 13.1.5 que $B$ est un groupe abéloïde sur $S$ au voisinage de $s$ si et seulement si sa fibre en $s$
\[
B_o=A_o/T_o
\]
est un groupe analytique abéloïde, de sorte que $A_o$ apparaît alors comme une extension d'un tel groupe $B_o$ par un tore $T_o$. \uline{Toute façon d'écrire} $A_o$ \uline{sous forme d'une extension d'un groupe abéloïde par un tore complexe} (*) \uline{donne donc naissance sur} $A$, \uline{au voisinage de} $s$, \uline{à une unique structure d'extension} (13.3.13) \uline{d'un groupe abéloïde} $B$ \uline{par un tore} $T$, \uline{induisant la structure d'extension donnée en} $s$.

13.4. \uline{Interprétation en termes de} ``\uline{structures de Hodge mixtes}'' \uline{de DELIGNE}.

Reprenons d'abord le dictionnaire 13.1.4, que nous allons reformuler d'une façon légèrement différente, mieux adaptée à l'introduction des ``structures de Hodge'' pertinentes à la situation. Pour ceci, soit $\mathbb R$ le faisceau des sections de $R$ sur $S$, et introduisons le $\underline O_S$-Module
\[
\text{(13.4.1)}\qquad \mathcal E(A)=\mathcal E=\underline R\otimes_{\mathbb Z}\underline O_S,
\]
qui est localement libre de type fini lorsque $R$ est localement constant. (NB Le faisceau $\mathcal E$ n'est pas nécessairement cohérent si $R$ n'est pas localement constant.) L'homomorphisme (13.1.3) $\underline R\longrightarrow\underline E$ définit alors

\margin{\unsure{493|1493}}
\margin{493}
\margin{89}
